% Blueprint content for the infinitary logic formalization

% Load generated node definitions from LeanArchitect
\input{../../.lake/build/blueprint/library/InfinitaryLogic}
\input{../../.lake/build/blueprint/library/InfinitaryLogicConditional}

\chapter{Infinitary Logic}

This project formalizes results in infinitary logic in Lean~4 using Mathlib.
The primary focus is the logic $\Lomegaone$, which extends first-order logic
with countable conjunctions and disjunctions, though the formalization
treats the more general $\Linf$ (allowing arbitrary conjunctions and
disjunctions) where possible.

\section{Main Results}

Headline results include:
\begin{itemize}
  \item \textbf{Scott characterization.}
    Every countable structure is characterized up to isomorphism by
    a sentence of $\Lomegaone$ (Theorem~\ref{thm:scott-characterizes}).
  \item \textbf{Scott rank below $\omegaone$.}
    The Scott rank of any countable structure is a countable ordinal
    (Theorem~\ref{thm:scottRank-lt-omega1}).
  \item \textbf{Karp's theorem.}
    Two structures are $\Linf$-equivalent if and only if there exists
    a back-and-forth system between them (Theorem~\ref{thm:karp-theorem}).
  \item \textbf{Model existence.}
    Every countable consistent set of sentences of $\Lomegaone$ in a countable
    language has a countable model (Theorem~\ref{thm:model-existence}).
  \item \textbf{Countable refinement hypothesis.}
    The back-and-forth hierarchy on any countable structure stabilizes
    in countably many steps (Theorem~\ref{thm:CRH}).
  \item \textbf{Silver's theorem and the Silver--Burgess dichotomy.}
    A Borel equivalence relation on a Polish space has countably many classes
    or a perfect set of pairwise-inequivalent points; on a standard Borel space
    the quotient has size $\leq\aleph_0$ or exactly $2^{\aleph_0}$
    (Theorem~\ref{thm:silver-burgess-proved}).
  \item \textbf{Morley counting.}
    An $\Lomegaone$ sentence has $\leq\aleph_1$ or exactly $2^{\aleph_0}$
    isomorphism classes of countable models (Theorem~\ref{thm:morley-counting}).
  \item \textbf{The Morley--Hanf theorem.}
    $\beth_{\omegaone}$ is a Hanf bound for every $\Lomegaone$ sentence,
    over an arbitrary language with no side hypotheses
    (Theorem~\ref{thm:morley-hanf}).
  \item \textbf{The exact Hanf number.}
    $\mathrm{Hanf}(\Lomegaone) = \beth_{\omegaone}$: the Morley--Hanf bound is
    sharp, by Marker's Exercise~5.3 beth ladder over the von Neumann hierarchy
    (Theorem~\ref{thm:hanf-exact}).
  \item \textbf{Small models of every infinite size.}
    A sentence with arbitrarily large models has, at every infinite $\kappa$, a
    model of size exactly $\kappa$ realizing only countably many complete
    $\Lomegaone$-types (Theorem~\ref{thm:small-models}).
  \item \textbf{Complete sentences and categoricity.}
    Every small model of $\varphi$ satisfies a complete sentence entailing
    $\varphi$; a $\kappa$-categorical $\varphi$ with arbitrarily large models
    has a complete completion with a $\kappa$-sized model, itself
    $\kappa$-categorical (Theorems~\ref{thm:complete-subclass},
    \ref{thm:categorical-complete}).
  \item \textbf{Craig interpolation.}
    An entailment between $\Lomegaone$ sentences over an \emph{arbitrary}
    language has an interpolant whose function and relation symbols each lie in
    the intersection of the two sides' occurrence sets
    (Theorem~\ref{thm:craig}).
  \item \textbf{Lyndon interpolation.}
    The polarity refinement, in the relation-polarity / logical-equality form:
    the interpolant's positively (negatively) occurring relation symbols lie in
    the intersection of the two sides' positive (negative) occurrences, with
    equality logical and unconstrained (Theorem~\ref{thm:lyndon}).
  \item \textbf{Boundedness and undefinability of well-ordering.}
    Chains of every countable length force a model with a relation-preserving
    map from $\mathbb{Q}$; hence the order types of well-ordered models are
    uniformly bounded below $\omegaone$, and no sentence defines the class of
    well-orders (Theorems~\ref{thm:wellordering-map},
    \ref{thm:wellordering-boundedness}, \ref{thm:wellordering-undefinable}).
  \item \textbf{The L\'opez--Escobar theorem.}
    Over a countable relational vocabulary, a class of coded countable
    structures is Borel and isomorphism-invariant if and only if it is the model
    class of a single $\Lomegaone$-sentence
    (Theorems~\ref{thm:lopez-escobar}, \ref{thm:lopez-escobar-iff}).
  \item \textbf{Non-Borelness of the countable well-order class.}
    The class of codes whose distinguished relation well-orders the carrier is
    not Borel in the logic space
    (Theorem~\ref{thm:wellordering-nonborel}).
\end{itemize}

\section{Back-and-Forth Equivalence}

Back-and-forth equivalence at ordinal levels, following Karp~\cite{karp1965}; the
infinitary languages themselves are developed in Karp~\cite{karp1964}.

\inputleannode{def:BFEquiv}

\section{Scott Formulas and Sentences}

Scott formulas and their characteristic property.  The theorem is
Scott's~\cite{scott1965}; the presentation here follows the modern exposition in
Marker~\cite[Ch.~2]{marker2016}.

\inputleannode{def:scottFormula}

\inputleannode{thm:scottFormula-iff}

\section{Self-Stabilization}

\inputleannode{thm:self-stabilization}

\section{Scott Sentence}

\inputleannode{def:stabilization-ordinal}

\inputleannode{def:scott-sentence}

\inputleannode{thm:scott-characterizes-of}

\chapter{Countable Refinement}

The countable refinement hypothesis and its consequences for Scott rank.

\section{Countable Refinement Hypothesis}

\inputleannode{thm:CRH}

\emph{Note:} Despite its name, the ``Countable Refinement Hypothesis'' is fully
proved in this formalization.  The name is historical: it was originally an
open conjecture about the countability of refinement ordinals.  The proof here
proceeds via self-stabilization and transfinite induction.

This countability result drives the downstream stabilization argument: since
only countably many ordinals can witness a strict refinement, the hierarchy
must stabilize below~$\omegaone$.

\inputleannode{thm:scott-characterizes}

\section{Scott Rank and Height}

\inputleannode{def:scottRank}

\inputleannode{thm:scottRank-lt-omega1}

\inputleannode{def:scottHeight}

\inputleannode{thm:scottHeight-lt-omega1}

\chapter{Karp's Theorem}

Karp's theorem characterizes $\Linf$-equivalence via back-and-forth
systems~\cite{karp1965}; the statement formalized here is Theorem~1.2.1 of
Keisler--Knight~\cite{keislerknight2004}.

The formalization states the equivalence at an arbitrary common branching carrier
equipped with codings of both structures, taking the disjoint sum as a corollary rather
than as the definition.  That is a presentation choice of this formalization, not part of
the cited statement: the mathematical content is Karp's.

\inputleannode{def:potential-iso}

\inputleannode{def:linf-equiv}

\inputleannode{thm:karp-theorem}

\chapter{Model Existence}

Model existence and completeness for $\Lomegaone$, following
Keisler~\cite{keisler1971}.

\inputleannode{thm:model-existence}

\inputleannode{thm:karp-completeness}

\section{Omitting Types}

\inputleannode{def:omits-type}

\inputleannode{thm:omitting-types}

\chapter{L\"owenheim--Skolem}

The downward L\"owenheim--Skolem theorem for $\Lomegaone$.

\inputleannode{def:naming-function-with-constants}

\inputleannode{thm:downward-ls}

\inputleannode{thm:downward-ls-theory}

\chapter{Admissible Fragments}

\emph{Scaffolding, not yet theorems.}  This chapter records the \emph{interfaces}
for compactness and completeness over countable admissible fragments of
$\Lomegaone$, following Barwise~\cite{barwise1975}; background on the Nadel bound
is in Nadel~\cite{nadel1974}.  The compactness principle is currently carried as a
\emph{structure field} and the Nadel bound as a \emph{typeclass hypothesis}, so the
nodes below state conditional interfaces rather than proved results.  Replacing
these by honest foundations is tracked in issues~\#18--\#20.

\inputleannode{def:admissible-fragment}

\inputleannode{thm:barwise-compactness}

\inputleannode{thm:barwise-completeness-ii}

\chapter{Hanf Numbers}

Hanf numbers measure how large a model a sentence must have before
it has arbitrarily large models.

\inputleannode{def:arb-large-models}

\inputleannode{def:hanf-bound}

\inputleannode{def:hanf-number}

\inputleannode{thm:hanf-existence}

\section{Morley--Hanf Bound}

\emph{Status: proved, unconditionally.}  The theorem \texttt{morley\_hanf}
(\ref{thm:morley-hanf}) states that $\beth_{\omegaone}$ is a Hanf bound for
every $\Lomegaone$ sentence, over an arbitrary language with no side
hypotheses.  No extraction is consumed at all: any injective
$\mathbb{N}$-sequence of the source is already fully indiscernible on the
Morley seed (\texttt{morleySeed\_indiscernibleOn}).  The single substantive
input is realizability of the seed's Ehrenfeucht--Mostowski tail-template
theory over every target order
(\texttt{morleySeedTailTemplateRealizable\_holds}).  The latter is a
Henkin-style construction in the spirit of Marker's presentation of Morley's
omitting-types argument: an $\omega$-stage completion of the countable
\emph{schema} sentence universe --- deciding each lifted template sentence and
pinning a disjunct for every positive countable disjunction and a refuted
conjunct for every negative countable conjunction --- whose quotient term model
carries a fully indiscernible sequence of Henkin constants satisfying a
restricted truth lemma.  Symbol countability is removed by running the
construction in the sublanguage generated by the sentence's own countably many
symbols and expanding the resulting model back to the full language.

The historical conditional forms are retained: the original
\texttt{morley\_hanf\_of\_transfer} (\ref{thm:morley-hanf-transfer}) with the
bundled transfer hypothesis, and the split bridge
\texttt{hasArbLargeModels\_of\_restricted\_extraction}
(\ref{thm:morley-hanf-restricted}) consuming the (since refuted-in-ZFC)
full-indiscernibility extraction residual plus a compactness oracle.

\inputleannode{thm:morley-hanf}

The milestones of the schema construction behind the seed-realizability input:

\inputleannode{thm:schema-truth-lemma}

\inputleannode{thm:schema-seq-indiscernible}

\inputleannode{thm:morley-seed-realizable}

\inputleannode{thm:morley-hanf-seed}

\inputleannode{thm:morley-hanf-transfer}

\inputleannode{thm:morley-hanf-restricted}

\section{Sharpness: the Exact Hanf Number}

\emph{Status: proved.}  The Morley--Hanf bound is exact:
$\mathrm{Hanf}(\Lomegaone) = \beth_{\omegaone}$
(\texttt{Lomega1omegaHanfNumber\_eq\_beth\_omega1}, \ref{thm:hanf-exact}).
The upper half (\ref{thm:hanf-upper-global}) is immediate from the Morley--Hanf
theorem.  The lower half (\ref{thm:hanf-lower-global}) is Marker's
Exercise~5.3: for each $\alpha < \omegaone$ a single $\Lomegaone$ sentence
over countably many constants, unary levels indexed by the well order
$\alpha+2$, and one binary extensional edge relation has models of size
exactly $\beth_{\alpha+1}$ at the maximum (\ref{thm:hanf-ladder-stage}).
The maximal model is the von Neumann hierarchy: the levels are interpreted as
the stages $V_{\omega+\beta}$, whose cardinalities Mathlib already computes,
shrunk to a $\mathrm{Type}$-$0$ carrier; the upper bound is a well-founded
induction on the level order with powerset steps at successors and countable
unions at limits.  The successor-cofinal supremum
$\sup_{\alpha<\omegaone}\beth_{\alpha+1} = \beth_{\omegaone}$ closes the
argument.

\inputleannode{thm:hanf-upper-global}

\inputleannode{thm:hanf-ladder-stage}

\inputleannode{thm:hanf-lower-global}

\inputleannode{thm:hanf-exact}

\section{Small Models}

\emph{Status: proved.}  Marker's Theorem~11.2
(\texttt{exists\_small\_model\_of\_hasArbLargeModels}, \ref{thm:small-models}): a sentence
with arbitrarily large models has models of every infinite size realizing only countably many
complete $\Lomegaone$-types.  The witnesses are the local Ehrenfeucht--Mostowski quotients of
the Morley-seed schema source over highly order-transitive skeletons (linear ordered fields,
supplied at every infinite cardinality by lexicographic Hahn-series subfields); the quotient is
equivariant under order automorphisms of the skeleton, tuples are classified by countably many
compressed term codes, and located term codes compute the carrier cardinality exactly.
Arbitrary languages reduce along a uniform collapsing hom.

\inputleannode{thm:small-models}

\chapter{Counting Models}

Counting the number of models of a Scott sentence up to isomorphism.
The theorem \texttt{morley\_counting\_dichotomy} (a counting-models
dichotomy in the style of Morley) is formalized in Lean but is not
included as a blueprint node; it serves as a schematic boundary result.

\emph{Status: schematic.}  The dichotomy statement quantifies over
cardinals and does not correspond to a single blueprint-trackable proof
obligation.

\inputleannode{thm:stabilization-bound-iso}

\inputleannode{thm:bounded-scott-height-iso}

\chapter{Proof System for Admissible Fragments}

A proof system for admissible fragments of $\Lomegaone$, following Barwise~\cite{barwise1975}.

\inputleannode{def:derivable}

\inputleannode{def:a-consistent}

\inputleannode{thm:proof-system-soundness}

\inputleannode{def:full-barwise-fragment}

\inputleannode{thm:consistency-property-full-fragment}

\inputleannode{thm:barwise-completeness-ii-syntactic}

\chapter{Descriptive Set Theory}

Borel complexity of satisfaction, back-and-forth equivalence, and isomorphism
on the coding space for countable structures.

\inputleannode{def:structure-space}

\inputleannode{thm:satisfaction-borel}

\inputleannode{thm:lopez-escobar-easy}

\section{L\'opez--Escobar}

\emph{Status: proved, both directions.}  The theorem is
L\'opez--Escobar's~\cite{lopezescobar1965}.  The hard direction follows Marker's
route (Theorem~4.25 of~\cite{marker2016}) rather than an induction on the Borel
hierarchy, which would not preserve isomorphism invariance.  An invariant Borel
class and its complement are analytic, hence branch projections of cylinder
trees along the query code; each tree is coded as a projective ($PC$) class over
the base language expanded by two \emph{disjoint tagged copies} of a functional
witness vocabulary, relationalized through the graph translation of the Craig
layer.  The base reducts of the coded class are exactly the original class ---
this converse is the sole consumer of isomorphism invariance --- and the two
tagged presentations have no common model, since gluing two expansions of a
shared base model and passing to a countable fragment-elementary substructure
would place a single code in a class and in its complement.  Craig separation
then produces a sentence of the shared vocabulary, which by the occurrence
calculus contains only graph images of base relation symbols and therefore
decodes back to an $\Lomegaone$-sentence over the base language.  Both final
inclusions use only the invariance-free forward presentation, so the whole
theorem's use of invariance stays concentrated in the disjointness lemma.

\inputleannode{thm:lopez-escobar}

\inputleannode{thm:lopez-escobar-iff}

\inputleannode{thm:lopez-escobar-action-iff}

\section{Complete Sentences and Categoricity}

\emph{Status: proved.}  The payoff of the small-model theorem (Marker, Theorem~11.2
applications): a small model's countable companion — built from realized-type isolators, one
countable controlling fragment, and the genuine fragment L\"owenheim--Skolem theorem — is
back-and-forth equivalent to it at every ordinal, and the companion's canonical Scott sentence
characterizes ARBITRARY models up to all-ordinal back-and-forth equivalence (the
arbitrary-target stabilization kernel, where the fragment machinery upgrades countable-target
Scott stabilization).  Hence every small model of $\varphi$ lies in a complete subclass of
$\mathrm{Mod}(\varphi)$, and $\kappa$-categorical sentences have $\kappa$-categorical
complete completions.

\inputleannode{thm:complete-subclass}

\inputleannode{thm:categorical-complete}

\inputleannode{thm:bfequiv-borel}

\inputleannode{thm:iso-borel}

\section{Standard Borel Structure}

The structure space carries the product topology, which makes it a Polish space
when the language has countably many relation symbols.

\inputleannode{thm:structure-space-polish}

The model class of any $\Lomegaone$ sentence inherits a standard Borel structure.

\inputleannode{thm:models-clopenable}

\inputleannode{thm:models-standard-borel}

\section{Counting Models}

The Silver--Burgess dichotomy, taken as a hypothesis, yields a counting theorem
for coded $\mathbb{N}$-models with bounded Scott height.

\inputleannode{def:silver-burgess-dichotomy}

\inputleannode{thm:silver-burgess-proved}

\inputleannode{def:structure-iso-setoid}

\inputleannode{def:iso-setoid}

\inputleannode{thm:counting-dichotomy}

\inputleannode{thm:finite-carrier-iso-borel}

\inputleannode{thm:counting-all-countable}

\inputleannode{thm:morley-counting}

The combined theorem extends the coded $\mathbb{N}$-model result to all countable
models by adding finite-carrier tiers (isomorphism on $\operatorname{Fin} n$
structures is Borel via the orbit of $\operatorname{Sym}(n)$).  Bridge theorems
(\texttt{codeModel}, \texttt{iso\_of\_codeModel\_eq}, \texttt{codeModel\_surjective})
establish that \texttt{AllCodedIsoClasses} faithfully represents all isomorphism
classes of countable models.  The Silver--Burgess dichotomy
(\texttt{silverBurgessDichotomy}) and Silver's theorem
(\texttt{gandy\_harrington\_for\_relation}) are \emph{fully proved}, via
Miller's classical category route (Kuratowski--Ulam, Mycielski, the
Kechris--Solecki--Todorcevic $G_0$-dichotomy); \texttt{morley\_counting} is
therefore unconditional, with axioms exactly
\texttt{[propext, Classical.choice, Quot.sound]}.

\section{Thinness and Perfect Sets}

The counting theorems above are stated cardinally.  The underlying dichotomy is
sharper: the ``many classes'' side is witnessed by a perfect set of pairwise
inequivalent points.  This section records the vocabulary for that witness, and
a criterion for ruling it out.

\inputleannode{def:perfect-antichain}

\inputleannode{def:cantor-antichain}

\inputleannode{def:thin-on}

\inputleannode{def:thin-nat-models}

The Cantor antichain is the load-bearing intermediary in both directions: it is
what the Cantor-scheme builders produce, and it is what a thinness proof must
refute.  The two chains are
\[
  \text{perfect antichain} \longleftrightarrow \text{Cantor antichain}
    \longrightarrow 2^{\aleph_0} \text{ classes},
\]
\[
  \mathtt{ThinRankAnalysis} \longrightarrow \text{no Cantor antichain}
    \longrightarrow \text{thinness},
\]
with $\mathtt{structureIsoSetoid}$ together with $\mathrm{ModelsOf}\ \varphi$
specializing the ambient notions to $\mathtt{IsThinOnNatModels}$.

\inputleannode{thm:cantor-to-perfect}

The two directions of that equivalence cost very different hypotheses.  Perfect
antichain $\to$ Cantor antichain goes through
$\mathtt{Perfect.exists\_nat\_bool\_injection}$ and needs a complete metric
space.  The converse needs only that the ambient space is Hausdorff, which is
why $\mathtt{HasCantorAntichainOn}$ is worth keeping as a separate notion: a
construction that produces one obtains the ambient perfect set without first
exhibiting a metric.

\inputleannode{def:thin-rank-analysis}

\inputleannode{thm:ranked-thinness}

\texttt{ThinRankAnalysis} packages sufficient evidence for thinness, and
\texttt{ThinRankAnalysis.isThinOn} proves the implication.  No concrete instance
of that package is supplied here.  Note that the cheap direction above does not
weaken those hypotheses: \texttt{isThinOn} consumes perfect $\to$ Cantor, the
complete-metric direction.

Cardinal equality alone does not establish thinness in ZFC.  Under $\mathsf{CH}$
we have $\aleph_1 = 2^{\aleph_0}$, so a quotient of exactly $\aleph_1$ classes is
compatible with a perfect antichain; under $\neg\mathsf{CH}$, exact $\aleph_1$
does rule one out.  Separately, the cardinal interface forgets the witness:
$\mathtt{SilverBurgessDichotomy}$ concludes with a cardinal disjunction and
retains nothing from which a perfect set could be extracted, so no theorem
parameterized by it can produce a perfect-set witness.

\section{Morley Counting with a Witness}

The way past that obstruction is to bypass the cardinal interface entirely and
apply the proved Polish-space core \texttt{silver\_core\_polish} directly.  The
model class is a \emph{Borel} subset of the structure space, hence not Polish as
a subtype, so a clopenable refinement is taken first, Silver runs there, and the
Cantor antichain is carried back to the ambient space --- the pipeline factored
as \texttt{silver\_countable\_or\_cantorAntichain}.  Only the Cantor form is
coarsened, where continuity is the sole topological clause; perfectness is
recovered ambiently afterwards, so nothing here asserts that perfectness
survives coarsening, which it does not.

\inputleannode{thm:morley-counting-coded-or-perfect}

\inputleannode{thm:counting-fin-models-or-perfect}

The two tiers differ only in the relation Silver is applied to.  At $\mathbb{N}$
it is back-and-forth equivalence at a level $\alpha < \omegaone$ --- what the
Scott stratification makes Borel --- and the antichain transfers to isomorphism,
which refines it.  At $\operatorname{Fin} n$ isomorphism is itself Borel, a
finite union of permutation graphs, so Silver applies to it directly.

\inputleannode{thm:morley-counting-or-perfect}

The finite-carrier alternative is not a formality: an infinite language can have
continuum-many $\operatorname{Fin} n$-models and no $\mathbb{N}$-models at all,
so a statement offering only an $\mathbb{N}$-tier perfect set would be false.
Since a perfect antichain forces continuum-many classes, the cardinal form of
Theorem~\ref{thm:morley-counting} is recovered as the corollary
\texttt{morley\_counting\_or\_perfect\_cardinal}, with the upper bound supplied
by the ambient space; the converse fails, which is the sense in which the
perfect alternatives carry strictly more information.  The parameterized
\texttt{morley\_counting} is retained unchanged.

\chapter{Craig Interpolation}

Craig interpolation for $\Lomegaone$: an entailment $r_1 \models r_2$ admits an
interpolant $\theta$ whose vocabulary lies in the symbols shared by the two
sides, with $r_1 \models \theta$ and $\theta \models r_2$.

The headline node holds over an \emph{arbitrary} language, with no hypotheses on
$L$: both the function and the relation symbols of $\theta$ lie in the
intersection of the two roots' occurrence sets.  It is the relational core
composed with a self-contained \emph{relationalization} layer: every $n$-ary
function symbol becomes an $(n{+}1)$-ary graph relation $G_f$, terms flatten to
existential graph formulas, the totality/functionality axioms of the mentioned
symbols translate the entailment (via structure reconstruction over the union of
the two function-symbol sets), and back-translation $G_f(\bar x,y) \mapsto
f(\bar x){=}y$ returns the interpolant to $L$ with the sharp occurrence bounds.
Its disjunction form is the PC-separation corollary
\texttt{craig\_pcSeparation}.

\inputleannode{thm:craig}

The relational core below is proved over a purely relational language with no
countability hypothesis, by an inseparability/finite-condition Henkin
construction (a constant-free interpolant strips through the empty-support root
gate), with ambient relation-symbol countability removed by passing to the
sublanguage generated by the two roots' countably many symbols.  Its disjunction
form, \texttt{craig\_pcSeparation\_relational}, is retained in exactly the
shared-vocabulary packaging consumed by the projective-class work.

\inputleannode{thm:craig-relational}

\section{Lyndon Interpolation}

\emph{Status: proved, arbitrary languages.}  The polarity refinement of Craig: an
interpolant may be required to respect the \emph{sign} of every relation
occurrence.  What is proved here is the \textbf{relation-polarity /
logical-equality form} of L\'opez--Escobar's Theorem~4.1: relation symbols
satisfy his clause~(.4) in full --- a symbol occurring positively (negatively) in
the interpolant occurs positively (negatively) in \emph{both} roots --- while
equality is treated as a \emph{logical} symbol, belonging to neither polarity
class, so his clause~(.3), which additionally constrains where equality may
occur, is deliberately not claimed.  Nor is any theory-level form claimed: for
$\Lomegaone$ that is false (L\'opez--Escobar's Theorem~6.3 exhibits disjoint
$PC_\Delta$ classes inseparable by an elementary class), so the sentence-level
statement is sharp.

Polarity is tracked by a signed occurrence traversal on the existing syntax ---
antecedents flip the sign, quantifiers and the countable connectives preserve it,
equality contributes nothing --- so no negation-normal form is constructed.  Its
semantic content is that growing the positively-occurring relations and shrinking
the negatively-occurring ones preserves truth.  The proof \emph{refines} the
Craig engine rather than repeating it: only the separator class changes, and the
countable-completion kernel is consumed unchanged.  The arbitrary-language step
reuses Craig's relationalization layer verbatim, the key point being that the
graph relations $G_f$, which may occur with either polarity, back-translate into
\emph{equalities} and so vanish from the polarity bookkeeping.

\inputleannode{thm:lyndon}

\inputleannode{thm:lyndon-relational}

\section{Malitz universal interpolation}

Malitz's Theorem~4.5 refines interpolation by \emph{quantifier class}: if the consequent is
universal, the interpolant may be taken universal too.  The engine is a side-labelled budgeted
pair --- Feferman's split-sequent invariant in semantic form --- in which the interpolant's class
is \emph{paid for} by the right root's quantifier budget rather than established afterwards.

The companion \emph{relative existential preservation} theorem (4.6) is not claimed; see the
successor issue.

\inputleannode{thm:malitz}

\chapter{Undefinability of Well-Ordering}

Marker's consistency-property construction (Theorem~4.26 and Corollary~4.27
of~\cite{marker2016}): an $\Lomegaone$-sentence whose models
contain well-ordered chains of every countable length has a model embedding the rationals
positively into the distinguished relation.  Consequently the well-ordered models of any
sentence have uniformly bounded order type, and no sentence of $\Lomegaone$ defines the
class of well-orders.

All three nodes hold over an \emph{arbitrary} language.  The engine is a dedicated
consistency property over the constants-expanded relational core, whose members carry the
positive rational diagram $\{d_q < d_r : q < r\}$ together with a finite remainder whose
mentioned rational constants sit, at every level $\alpha < \omega_1$, on a chain with
$\alpha$-sized margins in an approximating model (Marker's gap condition~$(*)$); the
fifteen closure fields are Exercise~4.28.  The fair Henkin enumeration and the quotient
term model shared with the Craig kernel realize the completed diagram, and $q \mapsto
d_q^M$ is the rational map.  Symbol countability is removed by the two-sorted generated
sublanguage, and function symbols by the Craig relationalization layer, under which the
distinguished base relation is preserved definitionally.

The three public results are deliberately separate in strength: the raw positive map (no
injectivity claimed — that is a corollary under irreflexivity), the uniform order-type
bound (the strong form), and undefinability (the weak form).  The stronger induced-copy
conclusion and the Borel-coded form of undefinability are tracked separately and are not
claimed by these nodes.

\inputleannode{thm:wellordering-map}

\inputleannode{thm:wellordering-boundedness}

\inputleannode{thm:wellordering-undefinable}

The Borel-coded form is the descriptive-set-theoretic payoff, and it needs
L\'opez--Escobar: the class of codes whose distinguished relation well-orders
$\mathbb{N}$ is not Borel.  Were it Borel it would be $\mathrm{ModelsOf}\
\varphi$ on codes; conjoining the infiniteness axiom (so no finite model
escapes) and seeding a countable fragment-elementary substructure with the
witnesses of any well-ordering defect --- two incomparable unequal elements, or
an infinite descending chain --- transports that defect to a code of the class,
so every model of the conjunction is in fact well-ordered.  Corollary~4.27 then
bounds all their order types by one countable $\alpha$, which the comparison
structure of type $\alpha + \omega$ exceeds.

\inputleannode{thm:wellordering-nonborel}

Boundedness itself survives without any definability hypothesis on the class.  A
sentence-defined reduct class is isomorphism-invariant, so an arbitrary analytic
$A$ can never be presented as $\mathrm{codeReduct}\ ''\ \mathrm{ModelsOf}\ \Theta$
exactly; but it can always be sandwiched between that class and an invariant
envelope $W \supseteq A$, and boundedness only ever consumes the sandwich.  Taking
$W$ to be the well-order class gives the classical boundedness theorem for
$\Sigma^1_1$ subsets of $\mathrm{WO}$.  Its immediate $A = \mathrm{WO}$ consequence
is that $\mathrm{WO}$ is not \emph{analytic} --- it would then bound its own order
types, which the supply of countably infinite order types exceeds --- and
non-Borelness follows because every Borel set is analytic.  The node above is
proved separately from L\'opez--Escobar and does not go through this one.

\inputleannode{thm:analytic-wellorder-boundedness}

\bibliographystyle{amsalpha}
\bibliography{refs}
